\subsection{Navigation Campaign Reproduction}\label{sec:nav-reproduction}

The navigation results are newly executed runs, separate from the archived physics
campaign. From the repository root, regenerate the four paired closed-loop cases, the common-truth
replay, thirty additional landing trials, fifteen ascent replays, and the navigation figures/tables with:
\begin{lstlisting}
.venv/bin/python -m pip install -e ".[dev]"
.venv/bin/python report_assets/generate_navigation.py
.venv/bin/python -m pytest tests/test_synthetic_eskf.py tests/sim/test_timing_contract.py tests/sim/test_multirate.py tests/sim/test_flight_calibration.py -q
\end{lstlisting}
The generator writes resolved configurations, seed manifests, full-rate navigation
histories, physics histories, and contact/error summaries under
\path{results/report_navigation/}. The navigation histories include native specific force and inertial rate at the navigation period.
The fixed-truth replay additionally saves the actual
IMU/GNSS/radar samples, all estimator states, bias estimates, and covariance diagonals.
\path{report_assets/navigation_metrics.json} records source SHA-256 hashes, the base Git
revision, numerical package versions, data hashes, and the evaluation conventions. These records identify the evaluated source and environment.
The figures and generated numerical tables in \path{report_assets/navigation_*} are
included in the report source tree; the \path{results/} histories can be regenerated with the same command
or supplied with the results archive to inspect the complete time histories.

For a normal standalone run with replay and sensor post-processing:
\begin{lstlisting}
.venv/bin/python scripts/run_single.py scenario=flight_calibrated_landing gnc/navigation=synthetic_eskf hydra.run.dir=results/navigation_landing
\end{lstlisting}
Replacing \texttt{synthetic\_eskf} with \texttt{identity} selects exact state feedback.
The standalone runner's additional sensor CSVs are generated after the flight; the ESKF
has already used its online sensor model during the closed-loop simulation.

The IMU rate and navigation period are explicit and must agree; the report settings use
100 Hz and 10 ms. Neither runner changes the period automatically. An independent
same-truth comparison and diagnostic plots can also be generated with:
\begin{lstlisting}
.venv/bin/python scripts/analyze_navigation.py scenario=project2_landing
\end{lstlisting}
This optional comparison uses the same physical sensor and rotating-frame propagation.
Its seeded replay is separate from the online samples that influenced a closed-loop run.
Scale-factor, bias-random-walk, outlier and gate settings are optional experiments; the
nominal report configuration keeps them disabled.
